% Die Titelseite
% Im folgenden kommen ein paar Variablen, die auszuf�llen sind
% Bisher steht dort nur Musterinhalt
% Au�erdem m�ssen zei Dateien erstellt werden, Bild/Logo/Emblem des Fachgebietes
% sowie der Universit�t

\newcommand{\trtitle}{Nachrichtenaustausch unter Angriff}
\newcommand{\trtype}{Diplomarbeit}
\newcommand{\trauthor}{Thomas Bosch}
\newcommand{\trstrasse}{Warschauer Str. 82}
\newcommand{\trmatrikelnummer}{197609}
\newcommand{\trort}{10243 Berlin}
\newcommand{\trbetreuerA}{Dipl.-Inform. Olaf Kroll-Peters}
\newcommand{\trbetreuerB}{Dipl.-Inform. Rainer Bye}
\newcommand{\trprof}{Prof.\ Dr.-Ing.\ habil.\ Sahin Albayrak}
\newcommand{\trfachgebiet}{Agententechnologien in betrieblichen Anwendungen und der Telekommunikation}
\newcommand{\trinstitut}{Wirtschaftsinformatik und Quantitative Methoden}
\newcommand{\trfakultaet}{IV Elektrotechnik und Informatik}
\newcommand{\truni}{Technische Universit�t Berlin}
\newcommand{\trdate}{\today}

\thispagestyle{empty}

% Kopfzeile mit Logos.
% Eventuell die \hspace{} je nach Logogr��e anpassen
\begin{tabular}{lcr}
  \includegraphics[scale=0.4]{tu_logo_rot} & % dein_unilogo.jpg/.eps im Verzeichnis "bilder" ablegen
  \hspace{3cm} \truni \hspace{2cm} &
  \includegraphics[scale=0.5]{dailogo} % dein_fglogo.jpg/.eps im Verzeichnis "bilder" ablegen, Fachgebietslogo
  \\
\end{tabular}

\rule{\textwidth}{0.4pt}

\vspace{2.5cm}
\begin{center}
  \textbf{\LARGE \trtitle}
\end{center}
\vspace{2cm}

\begin{center}
  \textbf{\trtype} \\
  am Fachgebiet \trfachgebiet \\
  \trprof \\
  Institut f�r \trinstitut \\
  Fakult�t \trfakultaet \\
  \truni \\[0.5cm]
  vorgelegt von \\
  \textbf{\trauthor}
\end{center}

\vspace{1cm}


\begin{center}
\begin{tabular}{ll}
Betreuer: & \trbetreuerA \\
 					& \trbetreuerB
\end{tabular}
\end{center}

\vfill

\begin{tabular}{l}
\trauthor \\
Matrikelnummer:  \trmatrikelnummer \\
\trstrasse \\
\trort
\end{tabular}

\rule{\textwidth}{0.4pt}
